\textbf{Proof.}
Fix \(e\in\{L,U\}\).  By def:moment-matrix, \(Q(c_e)\succeq0\), and ass:exposure-positivity makes \(\Pi^{-1}\) well defined.  Congruence therefore gives
\[
A_e=\Omega^{1/2}\Pi^{-1}Q(c_e)\Pi^{-1}\Omega^{1/2}\succeq0. \tag{1}
\]
Put \(Z_B=\Omega^{1/2}B\).  Since \(S_B=Z_B^{\top}Z_B\), cyclicity of trace and def:completion-pseudoinverse-loss give
\[
\begin{aligned}
n^2g_B^+(c_e)
&=\operatorname{tr}\!\left[(Z_B^{\top}Z_B)^+Z_B^{\top}A_eZ_B\right]\\
&=\operatorname{tr}\!\left[Z_B(Z_B^{\top}Z_B)^+Z_B^{\top}A_e\right]
=\operatorname{tr}(P_BA_e).                                      \tag{2}
\end{aligned}
\]
The singular-value decomposition of \(Z_B\) shows that
\[
P_B=Z_B(Z_B^{\top}Z_B)^+Z_B^{\top}
\]
is the orthogonal projector onto \(\operatorname{range}(Z_B)\).  It also shows
\[
\operatorname{rank}(P_B)=\operatorname{rank}(Z_B)
=\operatorname{rank}(Z_B^{\top}Z_B)
=\operatorname{rank}(S_B)\in\{0,1,2\}.                            \tag{3}
\]
The rank-two identification and equality of gains and losses are exactly lem:zhao-rank-two-stratum-identification.

It remains to prove the singular approximation without invoking cross-rank continuity.  Choose a block-separable representative \(B^{\circ}\) of a singular class and, using \(\mathcal B_Z\ne\varnothing\), choose a representative \(C^{\circ}\) of any class in \(\mathcal B_Z\).  For \(t\in\mathbb R\), set
\[
B_t^{\circ}=B^{\circ}+tC^{\circ}.
\]
This matrix remains block separable.  The determinant
\[
p(t)=\det\!\left[(Z_{B^{\circ}}+tZ_{C^{\circ}})^{\top}
(Z_{B^{\circ}}+tZ_{C^{\circ}})\right]
\]
is a polynomial whose leading coefficient is
\(\det(Z_{C^{\circ}}^{\top}Z_{C^{\circ}})>0\).  Hence it is not the zero polynomial and has only finitely many roots.  Choose \(t_j\downarrow0\) away from those roots.  Then \(B_{t_j}^{\circ}\) has rank-two covariance scale.  Independently normalize its two nonzero arm columns and pass to the ordinary quotient, obtaining \([B_j]\in\mathcal B_Z\).  This normalization leaves the range of \(Z_{B_{t_j}^{\circ}}\), and therefore its projector, unchanged.

The set of orthogonal rank-two projectors in finite dimension is compact, so after passing to and relabeling a subsequence, \(P_{B_j}\to P_{\infty}\).  If \(x=Z_{B^{\circ}}u\), then
\[
x_j=(Z_{B^{\circ}}+t_jZ_{C^{\circ}})u
\in\operatorname{range}(P_{B_j}),
\qquad x_j\longrightarrow x.
\]
Because an orthogonal projection gives the closest point in its range,
\[
\lVert P_{B_j}x-x\rVert_2
\le \lVert x_j-x\rVert_2\longrightarrow0.
\]
Thus \(P_{\infty}x=x\) for every \(x\in\operatorname{range}(Z_{B^{\circ}})\), so
\[
\operatorname{range}(P_B)\subseteq\operatorname{range}(P_{\infty}),
\qquad P_{\infty}-P_B\succeq0.                                   \tag{4}
\]
Combining (1), (2), and (4), simultaneously for \(e=L,U\), gives
\[
\lim_{j\to\infty}g_{B_j}(c_e)
=n^{-2}\operatorname{tr}(P_{\infty}A_e)
\ge n^{-2}\operatorname{tr}(P_BA_e)
=g_B^+(c_e).                                                       \tag{5}
\]
The uncontrolled term is independent of the basis.  Therefore
\[
\limsup_{j\to\infty}\ell_{B_j}(c_e)\le\ell_B^+(c_e)
\qquad(e=L,U).
\]
For two real sequences, the limit superior of their maximum is at most the maximum of their limit superiors.  Hence
\[
\limsup_{j\to\infty}w(B_j)
\le\max\{\ell_B^+(c_L),\ell_B^+(c_U)\}=w^+(B).
\]
This is a one-sided pointwise approximation; no lower semicontinuity at a rank drop has been assumed. \(\square\)